\newcommand{\h}{0mm}
\newcommand{\hh}{0mm}
\newcommand{\hhh}{0mm}

\newcommand{\figincremental}{
\begin{figure}[t]
\centering
\includegraphics[width=0.69\linewidth]{figures/incremental.pdf}\hfill%
\includegraphics[width=0.28\linewidth]{figures/incremental.jpg}
\caption{Our training starts with both the generator (\textsc{G}) and discriminator (\textsc{D}) having a low spatial resolution of 4$\times$4 pixels. As the training advances, we incrementally add layers to \textsc{G} and \textsc{D}, thus increasing the spatial resolution of the generated images. All existing layers remain trainable throughout the process. Here \framebox{\small $N\times N$} refers to convolutional layers operating on $N\times N$ spatial resolution. This allows stable synthesis in high resolutions and also speeds up training considerably. One the right we show six example images generated using progressive growing at $1024\times1024$.
}
\label{fig:incremental}
\end{figure}
}

\newcommand{\figfadein}{
\begin{figure}[t]
\centering
\includegraphics[width=0.7\linewidth]{figures/fadein.pdf}
\caption{When doubling the resolution of the generator (\textsc{G}) and discriminator (\textsc{D}) we fade in the new layers smoothly. This example illustrates the transition from $16\times16$ images (a) to $32\times32$ images (c). During the transition (b) we treat the layers that operate on the higher resolution like a residual block, whose weight $\alpha$ increases linearly from $0$ to $1$. Here \framebox{\small $2\times$} and \framebox{\small $0.5\times$} refer to doubling and halving the image resolution using nearest neighbor filtering and average pooling, respectively. The \framebox{\small toRGB} represents a layer that projects feature vectors to RGB colors and \framebox{\small fromRGB} does the reverse; both use $1\times1$ convolutions. 
When training the discriminator, we feed in real images that are downscaled to match the current resolution of the network. 
During a resolution transition, we interpolate between two resolutions of the real images, similarly to how the generator output combines two resolutions.
}
\label{fig:fadein}
\end{figure}
}

\newcommand{\figbedroomqualitycomparison}{
\begin{figure}
\adjincludegraphics[width=0.33\columnwidth,trim=0 0 {.25\width} {.25\height},clip]{figures/LSGAN-bedroom2.jpg}\hfill
\adjincludegraphics[width=0.33\columnwidth,trim=0 0 {.625\width} {.625\height},clip]{figures/WGAN-GP-bedroom.jpg}\hfill
\adjincludegraphics[width=0.33\columnwidth,trim=0 0 {0.0\width} {0\height},clip]{generated/figures/bedroom_curated.jpg}\\
\makebox[0.33\linewidth][c]{\footnotesize \cite{Mao2016} ($128\times128$)}\hfill
\makebox[0.33\linewidth][c]{\footnotesize \cite{Gulrajani2017} ($128\times128$)}\hfill
\makebox[0.33\linewidth][c]{\footnotesize Our ($256\times256$)}
\caption{Visual quality comparison in \textsc{LSUN bedroom}; pictures copied from the cited articles.}
\label{fig:bedroomqualitycomparison}
\end{figure}
}

\newcommand{\x}{}
\renewcommand{\x}{\hspace{1.8mm}}
\newcommand{\tabMNIST}{
\begin{table}[h]
\resizebox{\linewidth}{!}{
\begin{tabular}[H]{|l@{\hspace{1mm}}l|c|c||c|c|c|c|}\hline
\textbf{Arch} & \textbf{} & \textbf{GAN} & \textbf{+ unrolling} & \textbf{WGAN-GP} & \textbf{+ our norm} & \textbf{+ mb stddev} & \textbf{+ progression}\\\hline
\multirow{2}{*}{K/4} & \#  & $30.6\pm20.7$ & $372.2\pm20.7$ & $640.1\pm136.3$ & $856.7\pm50.4$ & $\mathbf{881.3\pm39.2}$ & $859.5\pm36.2$\\
                               & KL & $5.99\pm0.04$ & \x$4.66\pm0.46$ & $1.97\pm0.70$ & \x$1.10\pm$0.19 & \x$1.09\pm0.16$ & \x$\mathbf{1.05\pm0.09}$\\\hline
\multirow{2}{*}{K/2} & \#  & $628.0\pm140.9$ & $817.4\pm39.9$ & $772.4\pm146.5$ & $886.6\pm58.5$ & $918.3\pm30.2$ & $\mathbf{919.8\pm35.1}$\\
                               & KL & $2.58\pm0.75$ & \x$1.43\pm0.12$ & $1.35\pm0.55$ & \x$0.98\pm0.33$ & \x$0.89\pm0.21$ & \x$\mathbf{0.82\pm0.13}$\\\hline
\end{tabular}}
\caption{Results for MNIST discrete mode test using two tiny discriminators (K/4, K/2) defined by \protect{\cite{Metz2016}}. The number of covered modes (\#) and KL divergence from a uniform distribution are given as an average $\pm$ standard deviation over 8 random initializations. Higher is better for the number of modes, and lower is better for KL divergence.}
\label{tab:MNIST}
\end{table}
}

\newcommand{\figCIFARresults}{
\begin{figure}
\adjincludegraphics[width=1.0\columnwidth,trim=0 0 0 0,clip]{generated/figures/CIFAR10_unsupervised.jpg}
\caption{CIFAR10 images generated using a network that was trained unsupervised (no label conditioning), and achieves a record $8.80$ inception score. }
\label{fig:CIFARresults}
\end{figure}
}

\newcommand{\tabCIFARresults}{
\begin{table}[h]
\footnotesize
\makebox[0.49\columnwidth][c]{\textsc{Unsupervised}}\hfill
\makebox[0.49\columnwidth][c]{\textsc{Label conditioned}}\\
\begin{minipage}[c]{0.5275\linewidth}
\begin{tabular}{l@{\hspace{1mm}}l@{\hspace{0mm}}c}
\hline
Method & & Inception score\\\hline
ALI                                & \tiny{\citep{Dumoulin2016}}   & $5.34\pm0.05$\\	%
GMAN                            & \tiny{\citep{Durugkar2016}}  & $6.00\pm0.19$\\
Improved GAN              & \tiny{\citep{Salimans2016B}}  & $6.86\pm0.06$\\
CEGAN-Ent-VI              & \tiny{\citep{Dai2017}}             & $7.07\pm0.07$\\
LR-AGN                        & \tiny{\citep{Yang2017}}          & $7.17\pm0.17$\\
DFM                              & \tiny{\citep{Warde2017}}        & $7.72\pm0.13$\\
WGAN-GP                    & \tiny{\citep{Gulrajani2017}}    & $7.86\pm0.07$\\
Splitting GAN                & \tiny{\citep{Grinblat2017}}      & $7.90\pm0.09$\\\hline
\multicolumn{2}{l}{Our (best run)}                                  & $\mathbf{8.80\pm0.05}$\\
\multicolumn{2}{l}{Our (computed from 10 runs)}          & $\mathbf{8.56\pm0.06}$\\\hline
\end{tabular}
\end{minipage}
\hfill
\begin{minipage}[c]{0.5275\linewidth}
\begin{tabular}{l@{\hspace{1mm}}l@{\hspace{1mm}}c}
\hline
Method & & Inception score\\\hline
DCGAN                         & \tiny{\citep{Radford2015}}       & $6.58$\\			%
Improved GAN              & \tiny{\citep{Salimans2016B}}   & $8.09\pm0.07$\\
AC-GAN                        & \tiny{\citep{Odena2017}}        & $8.25\pm0.07$\\
SGAN                            & \tiny{\citep{Huang2016}}        & $8.59\pm0.12$\\
WGAN-GP	            & \tiny{\citep{Gulrajani2017}}    & $8.67\pm0.14$\\
Splitting GAN		    & \tiny{\citep{Grinblat2017}}      & $8.87\pm0.09$\\\hline
&&\\
&&\\
&&\\
&&\\
\end{tabular}
\end{minipage}\\
\caption{CIFAR10 inception scores, higher is better.}
\label{tab:CIFARresults}
\end{table}
}

\newcommand{\tabvariants}{
\begin{table}
\centering
% Autogenerated by generate_figures.py -- DO NOT EDIT DIRECTLY

\resizebox{\linewidth}{!}{%
\begin{tabular}{|l@{\hspace{1mm}}l|r@{\hspace{2.5mm}}r@{\hspace{2.5mm}}r@{\hspace{2.5mm}}r@{\hspace{2.5mm}}r@{\hspace{2.5mm}}|c|r@{\hspace{2.5mm}}r@{\hspace{2.5mm}}r@{\hspace{2.5mm}}r@{\hspace{2.5mm}}r@{\hspace{2.5mm}}|c|}
\hline
&
& \multicolumn{6}{c|}{\textsc{\textbf{CelebA}}}
& \multicolumn{6}{c|}{\textsc{\textbf{LSUN bedroom}}}
\\
\cline{3-14}
\multicolumn{2}{|l|}{\textbf{Training configuration}}
& \multicolumn{5}{c|}{\raisebox{0mm}[3.5mm][1mm]{Sliced Wasserstein distance $\times10^3$}}
& \multicolumn{1}{c|}{\raisebox{0mm}[3.5mm][1mm]{MS-SSIM}}
& \multicolumn{5}{c|}{\raisebox{0mm}[3.5mm][1mm]{Sliced Wasserstein distance $\times10^3$}}
& \multicolumn{1}{c|}{\raisebox{0mm}[3.5mm][1mm]{MS-SSIM}}
\\
&
& 128
& 64
& 32
& 16
& Avg
&  
& 128
& 64
& 32
& 16
& Avg
&  
\\
\hline
(a) & \cite{Gulrajani2017}
& 12.99
& 7.79
& 7.62
& 8.73
& 9.28
& 0.2854
& 11.97
& 10.51
& 8.03
& 14.48
& 11.25
& \textbf{0.0587}
\\
(b) & + Progressive growing
& 4.62
& \textbf{2.64}
& 3.78
& 6.06
& 4.28
& \textbf{0.2838}
& 7.09
& 6.27
& 7.40
& 9.64
& 7.60
& 0.0615
\\
(c) & + Small minibatch
& 75.42
& 41.33
& 41.62
& 26.57
& 46.23
& 0.4065
& 72.73
& 40.16
& 42.75
& 42.46
& 49.52
& 0.1061
\\
\hline
(d) & + Revised training parameters
& 9.20
& 6.53
& 4.71
& 11.84
& 8.07
& 0.3027
& 7.39
& 5.51
& 3.65
& 9.63
& 6.54
& 0.0662
\\
(e$^\ast$) & + Minibatch discrimination
& 10.76
& 6.28
& 6.04
& 16.29
& 9.84
& 0.3057
& 10.29
& 6.22
& 5.32
& 11.88
& 8.43
& 0.0648
\\
(e) & \hphantom{+} Minibatch stddev
& 13.94
& 5.67
& 2.82
& 5.71
& 7.04
& 0.2950
& 7.77
& 5.23
& 3.27
& 9.64
& 6.48
& 0.0671
\\
(f) & + Equalized learning rate
& 4.42
& 3.28
& 2.32
& 7.52
& 4.39
& 0.2902
& \textbf{3.61}
& 3.32
& \textbf{2.71}
& 6.44
& 4.02
& 0.0668
\\
(g) & + Pixelwise normalization
& \textbf{4.06}
& 3.04
& \textbf{2.02}
& \textbf{5.13}
& \textbf{3.56}
& 0.2845
& 3.89
& \textbf{3.05}
& 3.24
& \textbf{5.87}
& \textbf{4.01}
& 0.0640
\\
\hline
(h) & Converged
& 2.42
& 2.17
& 2.24
& 4.99
& 2.96
& 0.2828
& 3.47
& 2.60
& 2.30
& 4.87
& 3.31
& 0.0636
\\
\hline
\end{tabular}}

\caption{%
Sliced Wasserstein distance (SWD) %
between the generated and training images (Section~\ref{sec:metric}) and multi-scale structural similarity (MS-SSIM) among the generated images for several 
training setups at $128\times128$. For SWD, each column represents one level of the Laplacian pyramid, and the last one gives an average of the four distances. %
}
\label{tab:variants}
\end{table}
}

\newcommand{\figvariants}{
\begin{figure}
% Autogenerated by generate_figures.py -- DO NOT EDIT DIRECTLY

\adjincludegraphics[trim=0 0 0 0,clip,width=0.097\columnwidth]{generated/figures/variant-iwass.jpg}\hfill
\adjincludegraphics[trim=0 0 0 0,clip,width=0.097\columnwidth]{generated/figures/variant-lod.jpg}\hfill
\adjincludegraphics[trim=0 0 0 0,clip,width=0.097\columnwidth]{generated/figures/variant-mb16.jpg}\hfill
\adjincludegraphics[trim=0 0 0 0,clip,width=0.097\columnwidth]{generated/figures/variant-params.jpg}\hfill
\adjincludegraphics[trim=0 0 0 0,clip,width=0.097\columnwidth]{generated/figures/variant-mbdisc.jpg}\hfill
\adjincludegraphics[trim=0 0 0 0,clip,width=0.097\columnwidth]{generated/figures/variant-mbstat.jpg}\hfill
\adjincludegraphics[trim=0 0 0 0,clip,width=0.097\columnwidth]{generated/figures/variant-wscale.jpg}\hfill
\adjincludegraphics[trim=0 0 0 0,clip,width=0.097\columnwidth]{generated/figures/variant-pixelnorm.jpg}\hspace*{2mm}\hfill
\adjincludegraphics[trim=0 0 0 0,clip,width=0.146\columnwidth]{generated/figures/variant-nocrippling.jpg}\\\vspace{-1.0mm}
\makebox[0.097\columnwidth][c]{\scriptsize (a)}\hfill
\makebox[0.097\columnwidth][c]{\scriptsize (b)}\hfill
\makebox[0.097\columnwidth][c]{\scriptsize (c)}\hfill
\makebox[0.097\columnwidth][c]{\scriptsize (d)}\hfill
\makebox[0.097\columnwidth][c]{\scriptsize (e$\ast$)}\hfill
\makebox[0.097\columnwidth][c]{\scriptsize (e)}\hfill
\makebox[0.097\columnwidth][c]{\scriptsize (f)}\hfill
\makebox[0.097\columnwidth][c]{\scriptsize (g)}\hfill
\makebox[0.146\columnwidth][c]{\scriptsize (h) Converged}

\caption{%
(a) -- (g) \textsc{CelebA} examples corresponding to rows in Table~\protect\ref{tab:variants}. These are intentionally non-converged. (h) Our converged result.
Notice that some images show aliasing and some are not sharp -- this is a flaw of the dataset, which the model learns to replicate faithfully.
}
\label{fig:variants}
\end{figure}
}

\newcommand{\figvariantsAppendixCeleb}{
\begin{figure}
\ifarxivversion
	\input{generated/images_of_variants2_arxiv}
\else
    % Autogenerated by generate_figures.py -- DO NOT EDIT DIRECTLY

\rotatebox{90}{\makebox[25mm][c]{\tiny (a) \cite{Gulrajani2017}}}\hfill\adjincludegraphics[trim=0 0 0 0,clip,width=0.980\columnwidth]{generated/figures/variant-appendix-iwass.jpg}\\
\rotatebox{90}{\makebox[25mm][c]{\tiny (b) Progressive growing}}\hfill\adjincludegraphics[trim=0 0 0 0,clip,width=0.980\columnwidth]{generated/figures/variant-appendix-lod.jpg}\\
\rotatebox{90}{\makebox[25mm][c]{\tiny (c) Small minibatch}}\hfill\adjincludegraphics[trim=0 0 0 0,clip,width=0.980\columnwidth]{generated/figures/variant-appendix-mb16.jpg}\\
\rotatebox{90}{\makebox[25mm][c]{\tiny (d) Revised training parameters}}\hfill\adjincludegraphics[trim=0 0 0 0,clip,width=0.980\columnwidth]{generated/figures/variant-appendix-params.jpg}\\
\rotatebox{90}{\makebox[25mm][c]{\tiny (e$\ast$) Minibatch discrimination}}\hfill\adjincludegraphics[trim=0 0 0 0,clip,width=0.980\columnwidth]{generated/figures/variant-appendix-mbdisc.jpg}\\
\rotatebox{90}{\makebox[25mm][c]{\tiny (e) Minibatch stddev}}\hfill\adjincludegraphics[trim=0 0 0 0,clip,width=0.980\columnwidth]{generated/figures/variant-appendix-mbstat.jpg}\\
\rotatebox{90}{\makebox[25mm][c]{\tiny (f) Equalized learning rate}}\hfill\adjincludegraphics[trim=0 0 0 0,clip,width=0.980\columnwidth]{generated/figures/variant-appendix-wscale.jpg}\\
\rotatebox{90}{\makebox[25mm][c]{\tiny (g) Pixelwise normalization}}\hfill\adjincludegraphics[trim=0 0 0 0,clip,width=0.980\columnwidth]{generated/figures/variant-appendix-pixelnorm.jpg}\\

\fi
\vspace*{-3mm}
\caption{%
A larger set of generated images corresponding to the non-converged setups in Table~\protect\ref{tab:variants}.}
\label{fig:variants-appendix-celeb}
\end{figure}
}

\newcommand{\figplots}{
\begin{figure}
\hspace*{-1.5mm}%
\mbox{%
\includegraphics[width=46.8mm]{figures/plot-quality-fixed2.pdf}
\includegraphics[width=46.8mm]{figures/plot-quality-grow2.pdf}
\includegraphics[width=46.8mm]{figures/plot-perf-fixed-vs-grow2.pdf}}\vspace*{-.5mm}\\
\makebox[45mm][c]{\hspace*{5.5mm}\small(a)}\hfill
\makebox[45mm][c]{\hspace*{5.5mm}\small(b)}\hfill
\makebox[45mm][c]{\hspace*{4.5mm}\small(c)}
\caption{%
Effect of progressive growing on training speed and convergence. The timings were measured on a single-GPU setup using NVIDIA Tesla P100.
(a) Statistical similarity with respect to wall clock time for \cite{Gulrajani2017} using \textsc{CelebA} at $128\times128$ resolution. Each graph represents sliced Wasserstein distance on one level of the Laplacian pyramid, and the vertical line indicates the point where we stop the training in Table~\protect\ref{tab:variants}.
(b) Same graph with progressive growing enabled. The dashed vertical lines indicate points where we double the resolution of G and~D.
(c) Effect of progressive growing on the raw training speed in $1024\times1024$ resolution.
}
\label{fig:plots}
\end{figure}
}

\newcommand{\tabnetwork}{
\begin{table}
\centering
% Autogenerated by generate_figures.py -- DO NOT EDIT DIRECTLY

\scriptsize
\begin{tabular}{|lccr|}
\hline
\textbf{Generator} & Act. & Output shape & Params \\
\hline
Latent vector & -- & $\makebox[\widthof{512}][c]{512}\times\makebox[\widthof{1024}][c]{1}\times\makebox[\widthof{1024}][c]{1}$ & -- \\
Conv $4\times4$ & {\tiny LReLU} & $\makebox[\widthof{512}][c]{512}\times\makebox[\widthof{1024}][c]{4}\times\makebox[\widthof{1024}][c]{4}$ & 4.2M \\
Conv $3\times3$ & {\tiny LReLU} & $\makebox[\widthof{512}][c]{512}\times\makebox[\widthof{1024}][c]{4}\times\makebox[\widthof{1024}][c]{4}$ & 2.4M \\
\hline
Upsample & -- & $\makebox[\widthof{512}][c]{512}\times\makebox[\widthof{1024}][c]{8}\times\makebox[\widthof{1024}][c]{8}$ & -- \\
Conv $3\times3$ & {\tiny LReLU} & $\makebox[\widthof{512}][c]{512}\times\makebox[\widthof{1024}][c]{8}\times\makebox[\widthof{1024}][c]{8}$ & 2.4M \\
Conv $3\times3$ & {\tiny LReLU} & $\makebox[\widthof{512}][c]{512}\times\makebox[\widthof{1024}][c]{8}\times\makebox[\widthof{1024}][c]{8}$ & 2.4M \\
\hline
Upsample & -- & $\makebox[\widthof{512}][c]{512}\times\makebox[\widthof{1024}][c]{16}\times\makebox[\widthof{1024}][c]{16}$ & -- \\
Conv $3\times3$ & {\tiny LReLU} & $\makebox[\widthof{512}][c]{512}\times\makebox[\widthof{1024}][c]{16}\times\makebox[\widthof{1024}][c]{16}$ & 2.4M \\
Conv $3\times3$ & {\tiny LReLU} & $\makebox[\widthof{512}][c]{512}\times\makebox[\widthof{1024}][c]{16}\times\makebox[\widthof{1024}][c]{16}$ & 2.4M \\
\hline
Upsample & -- & $\makebox[\widthof{512}][c]{512}\times\makebox[\widthof{1024}][c]{32}\times\makebox[\widthof{1024}][c]{32}$ & -- \\
Conv $3\times3$ & {\tiny LReLU} & $\makebox[\widthof{512}][c]{512}\times\makebox[\widthof{1024}][c]{32}\times\makebox[\widthof{1024}][c]{32}$ & 2.4M \\
Conv $3\times3$ & {\tiny LReLU} & $\makebox[\widthof{512}][c]{512}\times\makebox[\widthof{1024}][c]{32}\times\makebox[\widthof{1024}][c]{32}$ & 2.4M \\
\hline
Upsample & -- & $\makebox[\widthof{512}][c]{512}\times\makebox[\widthof{1024}][c]{64}\times\makebox[\widthof{1024}][c]{64}$ & -- \\
Conv $3\times3$ & {\tiny LReLU} & $\makebox[\widthof{512}][c]{256}\times\makebox[\widthof{1024}][c]{64}\times\makebox[\widthof{1024}][c]{64}$ & 1.2M \\
Conv $3\times3$ & {\tiny LReLU} & $\makebox[\widthof{512}][c]{256}\times\makebox[\widthof{1024}][c]{64}\times\makebox[\widthof{1024}][c]{64}$ & 590k \\
\hline
Upsample & -- & $\makebox[\widthof{512}][c]{256}\times\makebox[\widthof{1024}][c]{128}\times\makebox[\widthof{1024}][c]{128}$ & -- \\
Conv $3\times3$ & {\tiny LReLU} & $\makebox[\widthof{512}][c]{128}\times\makebox[\widthof{1024}][c]{128}\times\makebox[\widthof{1024}][c]{128}$ & 295k \\
Conv $3\times3$ & {\tiny LReLU} & $\makebox[\widthof{512}][c]{128}\times\makebox[\widthof{1024}][c]{128}\times\makebox[\widthof{1024}][c]{128}$ & 148k \\
\hline
Upsample & -- & $\makebox[\widthof{512}][c]{128}\times\makebox[\widthof{1024}][c]{256}\times\makebox[\widthof{1024}][c]{256}$ & -- \\
Conv $3\times3$ & {\tiny LReLU} & $\makebox[\widthof{512}][c]{64}\times\makebox[\widthof{1024}][c]{256}\times\makebox[\widthof{1024}][c]{256}$ & 74k \\
Conv $3\times3$ & {\tiny LReLU} & $\makebox[\widthof{512}][c]{64}\times\makebox[\widthof{1024}][c]{256}\times\makebox[\widthof{1024}][c]{256}$ & 37k \\
\hline
Upsample & -- & $\makebox[\widthof{512}][c]{64}\times\makebox[\widthof{1024}][c]{512}\times\makebox[\widthof{1024}][c]{512}$ & -- \\
Conv $3\times3$ & {\tiny LReLU} & $\makebox[\widthof{512}][c]{32}\times\makebox[\widthof{1024}][c]{512}\times\makebox[\widthof{1024}][c]{512}$ & 18k \\
Conv $3\times3$ & {\tiny LReLU} & $\makebox[\widthof{512}][c]{32}\times\makebox[\widthof{1024}][c]{512}\times\makebox[\widthof{1024}][c]{512}$ & 9.2k \\
\hline
Upsample & -- & $\makebox[\widthof{512}][c]{32}\times\makebox[\widthof{1024}][c]{1024}\times\makebox[\widthof{1024}][c]{1024}$ & -- \\
Conv $3\times3$ & {\tiny LReLU} & $\makebox[\widthof{512}][c]{16}\times\makebox[\widthof{1024}][c]{1024}\times\makebox[\widthof{1024}][c]{1024}$ & 4.6k \\
Conv $3\times3$ & {\tiny LReLU} & $\makebox[\widthof{512}][c]{16}\times\makebox[\widthof{1024}][c]{1024}\times\makebox[\widthof{1024}][c]{1024}$ & 2.3k \\
Conv $1\times1$ & linear & $\makebox[\widthof{512}][c]{3}\times\makebox[\widthof{1024}][c]{1024}\times\makebox[\widthof{1024}][c]{1024}$ & 51 \\
\hline
\multicolumn{3}{|l}{Total trainable parameters} & \textbf{23.1M} \\
\hline
\multicolumn{4}{l}{} \\
\multicolumn{4}{l}{} \\
\end{tabular}
\hfill
\begin{tabular}{|lccr|}
\hline
\textbf{Discriminator} & Act. & Output shape & Params \\
\hline
Input image & -- & $\makebox[\widthof{512}][c]{3}\times\makebox[\widthof{1024}][c]{1024}\times\makebox[\widthof{1024}][c]{1024}$ & -- \\
Conv $1\times1$ & {\tiny LReLU} & $\makebox[\widthof{512}][c]{16}\times\makebox[\widthof{1024}][c]{1024}\times\makebox[\widthof{1024}][c]{1024}$ & 64 \\
Conv $3\times3$ & {\tiny LReLU} & $\makebox[\widthof{512}][c]{16}\times\makebox[\widthof{1024}][c]{1024}\times\makebox[\widthof{1024}][c]{1024}$ & 2.3k \\
Conv $3\times3$ & {\tiny LReLU} & $\makebox[\widthof{512}][c]{32}\times\makebox[\widthof{1024}][c]{1024}\times\makebox[\widthof{1024}][c]{1024}$ & 4.6k \\
Downsample & -- & $\makebox[\widthof{512}][c]{32}\times\makebox[\widthof{1024}][c]{512}\times\makebox[\widthof{1024}][c]{512}$ & -- \\
\hline
Conv $3\times3$ & {\tiny LReLU} & $\makebox[\widthof{512}][c]{32}\times\makebox[\widthof{1024}][c]{512}\times\makebox[\widthof{1024}][c]{512}$ & 9.2k \\
Conv $3\times3$ & {\tiny LReLU} & $\makebox[\widthof{512}][c]{64}\times\makebox[\widthof{1024}][c]{512}\times\makebox[\widthof{1024}][c]{512}$ & 18k \\
Downsample & -- & $\makebox[\widthof{512}][c]{64}\times\makebox[\widthof{1024}][c]{256}\times\makebox[\widthof{1024}][c]{256}$ & -- \\
\hline
Conv $3\times3$ & {\tiny LReLU} & $\makebox[\widthof{512}][c]{64}\times\makebox[\widthof{1024}][c]{256}\times\makebox[\widthof{1024}][c]{256}$ & 37k \\
Conv $3\times3$ & {\tiny LReLU} & $\makebox[\widthof{512}][c]{128}\times\makebox[\widthof{1024}][c]{256}\times\makebox[\widthof{1024}][c]{256}$ & 74k \\
Downsample & -- & $\makebox[\widthof{512}][c]{128}\times\makebox[\widthof{1024}][c]{128}\times\makebox[\widthof{1024}][c]{128}$ & -- \\
\hline
Conv $3\times3$ & {\tiny LReLU} & $\makebox[\widthof{512}][c]{128}\times\makebox[\widthof{1024}][c]{128}\times\makebox[\widthof{1024}][c]{128}$ & 148k \\
Conv $3\times3$ & {\tiny LReLU} & $\makebox[\widthof{512}][c]{256}\times\makebox[\widthof{1024}][c]{128}\times\makebox[\widthof{1024}][c]{128}$ & 295k \\
Downsample & -- & $\makebox[\widthof{512}][c]{256}\times\makebox[\widthof{1024}][c]{64}\times\makebox[\widthof{1024}][c]{64}$ & -- \\
\hline
Conv $3\times3$ & {\tiny LReLU} & $\makebox[\widthof{512}][c]{256}\times\makebox[\widthof{1024}][c]{64}\times\makebox[\widthof{1024}][c]{64}$ & 590k \\
Conv $3\times3$ & {\tiny LReLU} & $\makebox[\widthof{512}][c]{512}\times\makebox[\widthof{1024}][c]{64}\times\makebox[\widthof{1024}][c]{64}$ & 1.2M \\
Downsample & -- & $\makebox[\widthof{512}][c]{512}\times\makebox[\widthof{1024}][c]{32}\times\makebox[\widthof{1024}][c]{32}$ & -- \\
\hline
Conv $3\times3$ & {\tiny LReLU} & $\makebox[\widthof{512}][c]{512}\times\makebox[\widthof{1024}][c]{32}\times\makebox[\widthof{1024}][c]{32}$ & 2.4M \\
Conv $3\times3$ & {\tiny LReLU} & $\makebox[\widthof{512}][c]{512}\times\makebox[\widthof{1024}][c]{32}\times\makebox[\widthof{1024}][c]{32}$ & 2.4M \\
Downsample & -- & $\makebox[\widthof{512}][c]{512}\times\makebox[\widthof{1024}][c]{16}\times\makebox[\widthof{1024}][c]{16}$ & -- \\
\hline
Conv $3\times3$ & {\tiny LReLU} & $\makebox[\widthof{512}][c]{512}\times\makebox[\widthof{1024}][c]{16}\times\makebox[\widthof{1024}][c]{16}$ & 2.4M \\
Conv $3\times3$ & {\tiny LReLU} & $\makebox[\widthof{512}][c]{512}\times\makebox[\widthof{1024}][c]{16}\times\makebox[\widthof{1024}][c]{16}$ & 2.4M \\
Downsample & -- & $\makebox[\widthof{512}][c]{512}\times\makebox[\widthof{1024}][c]{8}\times\makebox[\widthof{1024}][c]{8}$ & -- \\
\hline
Conv $3\times3$ & {\tiny LReLU} & $\makebox[\widthof{512}][c]{512}\times\makebox[\widthof{1024}][c]{8}\times\makebox[\widthof{1024}][c]{8}$ & 2.4M \\
Conv $3\times3$ & {\tiny LReLU} & $\makebox[\widthof{512}][c]{512}\times\makebox[\widthof{1024}][c]{8}\times\makebox[\widthof{1024}][c]{8}$ & 2.4M \\
Downsample & -- & $\makebox[\widthof{512}][c]{512}\times\makebox[\widthof{1024}][c]{4}\times\makebox[\widthof{1024}][c]{4}$ & -- \\
\hline
Minibatch stddev & -- & $\makebox[\widthof{512}][c]{513}\times\makebox[\widthof{1024}][c]{4}\times\makebox[\widthof{1024}][c]{4}$ & -- \\
Conv $3\times3$ & {\tiny LReLU} & $\makebox[\widthof{512}][c]{512}\times\makebox[\widthof{1024}][c]{4}\times\makebox[\widthof{1024}][c]{4}$ & 2.4M \\
Conv $4\times4$ & {\tiny LReLU} & $\makebox[\widthof{512}][c]{512}\times\makebox[\widthof{1024}][c]{1}\times\makebox[\widthof{1024}][c]{1}$ & 4.2M \\
Fully-connected & linear & $\makebox[\widthof{512}][c]{1}\times\makebox[\widthof{1024}][c]{1}\times\makebox[\widthof{1024}][c]{1}$ & 513 \\
\hline
\multicolumn{3}{|l}{Total trainable parameters} & \textbf{23.1M} \\
\hline
\end{tabular}

\caption{%
Generator and discriminator that we use with \textsc{CelebA-HQ} to generate $1024\times1024$ images.
}
\label{tab:network}
\end{table}
}

\newcommand{\figcelebprocess}{
\begin{figure}
\includegraphics[width=1.0\linewidth]{figures/process-all.pdf}\\
\makebox[22.8mm][c]{(a)}\hfill
\makebox[11.2mm][c]{(b)}\hfill
\makebox[23.0mm][c]{(c)}\hfill
\makebox[23.0mm][c]{(d)}\hfill
\makebox[23.0mm][c]{(e)}\hfill
\makebox[34.0mm][c]{(f)}
\caption{%
Creating the \textsc{CelebA-HQ} dataset.
We start with a JPEG image (a) from the CelebA in-the-wild dataset.
We improve the visual quality (b,top) through JPEG artifact removal (b,middle) and 4x super-resolution (b,bottom).
We then extend the image through mirror padding (c) and Gaussian filtering (d) to produce a visually pleasing depth-of-field effect.
Finally, we use the facial landmark locations to select an appropriate crop region (e) and perform high-quality resampling to obtain the final image at $1024\times1024$ resolution (f).
}
\label{fig:celeb-process}
\end{figure}
}

\newcommand{\figcelebHQResults}{
\begin{figure}
\centering
\includegraphics[width=1.0\linewidth]{generated/figures/celeb1024-results.jpg}
\caption{$1024\times1024$ images generated using the \textsc{CelebA-HQ} dataset. See Appendix~\ref{app:CelebHQ} for a larger set of results, and the accompanying video for latent space interpolations. }
\label{fig:celebHQResults}
\end{figure}
}

\newcommand{\figcelebHQAppendix}{
\begin{figure}
\centering
\ifarxivversion
	\placeholderfig{1.0\linewidth}{210mm}{\arxivDisclaimer}\\
\else
	\includegraphics[width=1.0\linewidth]{generated/figures/celeb1024-appendix.jpg}
\fi
\caption{Additional $1024\times1024$ images generated using the \textsc{CelebA-HQ} dataset. 
Sliced Wasserstein Distance (SWD) $\times10^3$ for levels 1024, \ldots, 16: 7.48, 7.24, 6.08, 3.51, 3.55, 3.02, 7.22, for which the average is 5.44. 
Fr\'echet Inception Distance (FID) computed from 50K images was 7.30.
See the video for latent space interpolations.
}
\label{fig:celebHQAppendxi}
\end{figure}
}

\newcommand{\figcelebHQNN}{
\begin{figure}
\centering
\includegraphics[width=1.0\linewidth]{generated/figures/NN2b-top.jpg}\\
\vspace*{1.5mm}
\includegraphics[width=1.0\linewidth]{generated/figures/NN2b-bottom.jpg}
\caption{Top: Our \textsc{CelebA-HQ} results. \FINAL{Next five rows: Nearest neighbors found from the training data, based on feature-space distance. We used activations from five VGG layers, as suggested by \cite{Chen2017}.
Only the crop highlighted in bottom right image was used for comparison in order to exclude image background and focus the search on matching facial features.}}
\label{fig:celebHQNN}
\end{figure}
}

